%==============================================================================
\section{Discussion}\label{sec:disc}
%==============================================================================

\subsection{Bound-limited, not implementation-limited}\label{ssec:disc-bound}

A method can fail to solve an instance for two reasons. It can fail to find a good
solution, or it can find one and fail to prove that no better one exists. The
distinction matters because the two failures call for different remedies. The first
is answered by better search: heuristics, warm starts, branching rules. The second is
answered only by a better bound.

For the SSP the campaign answers this question unambiguously, and it answers it four
times over.

The first answer is in the direct measurement. On the runs that hit the time limit and
whose optimum is known from elsewhere, the Benders solver's incumbent already equals
that optimum $66\%$ of the time across its nine configurations, and $59\%$ of the time
for the plain one. In those runs the solver has already found the answer and spends the
remaining hour failing to certify it.

The second answer is in the ablation. Four strengthenings were implemented. The one
that improves the incumbent, a hybrid genetic search run at the root, gains thirty-one
instances. The three that aim at the dual bound --- fractional cuts, Pareto-optimal
lifting, a conflict-graph row --- are neutral or harmful, and the worst of them costs
thirty-seven. A solver whose difficulty was finding good solutions would not behave
this way.

The third answer is the split by coverage-bound tightness. The same solver closes
$94\%$ of the instances on which the elementary bound $q$ happens to equal the
optimum, and $40\%$ of those on which it does not. The classification is by a single structural property rather than by size or density,
and it is the property Section~\ref{sec:gap} identifies as governing the problem.

The fourth answer comes from the other method entirely. The branch-and-price prototype
reports its root relaxation on every run, and on all $2{,}114$ runs of the
position--configuration formulation that value is exactly $q$. A different algorithm,
a different formulation and a different solver arrive at the same number.

The natural objection is that this is a statement about one implementation. It is not,
and Proposition~\ref{prop:pcf} is why. The position--configuration relaxation, arrived
at from a completely different direction, equals $q$ exactly --- proved in
Section~\ref{ssec:pos} and then measured on the benchmark. So does the multicommodity
relaxation of \citet{daSilva2024}, proved by its own authors. Three independent routes reach the same value, which makes it a property of the
problem rather than of any one implementation.
A configuration-based exact method for the SSP is limited by the strength of its bound
and not by its implementation.

The fractional-cut result is the strongest form of this conclusion. Separating at
fractional points is the standard remedy when a Benders scheme tails off, and it does
what it is supposed to do locally: sixty-seven million cuts reduce the node count by a
factor of twenty. It nonetheless raises the bound at termination on no instance at all.
The remedy is not being applied badly; it is being applied to the wrong obstacle.

\subsection{Where a stronger bound could come from}\label{ssec:disc-locality}

Section~\ref{ssec:disc-bound} establishes that the methods of this report are limited
by their bound. That is a diagnosis of the symptom. This subsection asks what causes
it, and measures what a different family of inequalities could be worth. The answer
turns out to be structural rather than incidental, which is why no amount of tuning
moved it.

\paragraph{The cost is not a local quantity.}
Write $\beta_t$ for the number of maximal blocks of consecutive positions during which
tool $t$ sits in the magazine. Every insertion of $t$ begins one such block, so
\[
  \Zst \;=\; \sum_{t\in\U}\beta_t ,
\qquad\text{and}\qquad
  \beta_t\ge 1 \;\text{ for every used tool}
\]
recovers the coverage bound exactly. Everything above $q$ is therefore a
\emph{re-insertion}: a tool evicted and brought back. A re-insertion happens when the
jobs needing a tool are pushed apart by other jobs whose requirements cannot sit
alongside it. That is a property of the global interleaving of the schedule. It is not
a property of which pairs of jobs happen to be adjacent.

\paragraph{Every inequality in this report is local.}
Against that, consider what the two methods actually give their relaxations.

The Benders master carries one row $\theta\ge\sum_{ij}w_{ij}x_{ij}$, which is indexed
by arcs; the triplet rows $\theta\ge w_{ijk}(x_{ij}+x_{jk}-1)$, which are the weak
linearisation of a product and are worth nothing at a fractional point where both arcs
sit at one half; and the constant coverage row. Every optimality cut it generates is
also arc-supported, with coefficient $\sum_t\lambda_{ijt}$ on $x_{ij}$. So the master
learns only adjacency-indexed facts about a cost that is not adjacency-indexed.

The position--configuration relaxation prices one configuration per position and
couples positions only through per-tool presence rows, so its linear relaxation mixes
configurations fractionally at each position and pays no transition. That is why
Proposition~\ref{prop:pcf} holds, and it is why the measured root value is $q$ on every
one of $2{,}114$ runs.

PTF is the one construction here that ever exceeds $q$, and the reason is the same one
stated positively: a column carries a transition, so the relaxation cannot avoid paying
for it. It exceeds $q$ rarely because the chaining between consecutive pairs is enforced
tool by tool rather than configuration by configuration, which restores the fractional
mixing one level up.

\begin{observation}[The coverage row is carrying the master alone]\label{obs:rootq}
On $9{,}084$ of the $10{,}396$ Benders runs for which a root relaxation value was
recorded --- $87.4\%$ --- that value equals $\lvert\U\rvert$ exactly. The pairwise row
and the triplet rows contribute nothing at the root on almost nine runs in ten. Where
the pairwise structure does bite it bites hard, by $7.70$ on average, but that is
confined to the dense instances.
\end{observation}

\paragraph{A family that is not local.}
Let $W=[p,r]$ be a contiguous block of positions with $p\ge2$. Write $w_{t,k}$ for the
insertion of tool $t$ at position $k$ and $a_{t,k}$ for its presence, as in the
position-indexed formulation of Section~\ref{ssec:pos}.

\begin{proposition}[Window inequality]\label{prop:window}
For every contiguous block of positions $W=[p,r]$ and every feasible point of PCF,
\begin{equation}\label{eq:window}
  \sum_{t\in\U}\ \sum_{k=p}^{r} w_{t,k}
  \;\;\ge\;\;
  \bigl\lvert\{t: a_{t,k}=1 \text{ for some } k\in W\}\bigr\rvert
  \;-\; \sum_{t\in\U} a_{t,p-1}.
\end{equation}
\end{proposition}

\begin{proof}
Let $S$ be the set of tools present at some position of $W$ and $A$ the set present at
position $p-1$. Fix $t\in S\setminus A$. Then $a_{t,p-1}=0$ and $a_{t,k}=1$ for some
$k\in W$, so there is a smallest such $k$, and at it $a_{t,k}-a_{t,k-1}=1$. Constraint
(W) gives $w_{t,k}\ge1$. Distinct tools contribute to distinct terms, so the left-hand
side is at least $\lvert S\setminus A\rvert\ge\lvert S\rvert-\lvert A\rvert$, which is
the right-hand side.
\end{proof}

The inequality linearises with one auxiliary variable per tool and window. Two
properties make the family useful. It constrains a stretch of the schedule rather than
a single point, so it can charge a re-insertion that no arc-indexed or single-position
row can see. And it contains only master variables, so it is a \emph{robust} cut in the
sense of \citet{Poggi2011}: its dual never reaches the pricing problem, and the
set-union knapsack oracle of Section~\ref{ssec:pos} is unchanged. The per-tool counting
rows (T) already in the master are the case $W=[2,n]$ summed over tools, so
Proposition~\ref{prop:window} generalises them.

The implementation was additionally checked against the loading rule at $1{,}200{,}816$
integer points --- every window of every sequence of a bank of random instances --- with
no violation. That is a test of the code, not of the proposition.

\paragraph{What it would be worth.}
A family is worth implementing only if the bound it could deliver is materially better
than the bound already available. That can be settled before writing any solver code,
by computing for each family the best bound it could possibly give and comparing it
against the optimum. On an instance small enough to enumerate, the ceiling of a family
is the minimum over sequences of the strongest bound that family certifies for that
sequence: for the arc-supported family this is the pairwise weight solved to optimality
as a Hamiltonian path, and for the window family it is the best decomposition of the
sequence into disjoint windows, an $O(n^2)$ dynamic program.
Table~\ref{tab:ceilings} reports both on the $62$ instances of the Laporte3 family whose
coverage bound is loose.

\begin{table}[h]\centering
\footnotesize
\begin{tabular}{@{}lrrrr@{}}
\toprule
Family & mean & median & best case & no gain on\\
\midrule
arc-supported (what the solvers carry today) & $8.9\%$  & $0\%$  & $60.0\%$ & $46/62$\\
window                                       & $30.9\%$ & $40\%$ & $66.7\%$ & $18/62$\\
\bottomrule
\end{tabular}
\caption{The fraction of the gap $\Zst-q$ that each family could close at best, on the
$62$ Laporte3 instances with $n=8$ whose coverage bound is loose. In absolute terms the
arc family accounts for $54$ of the $305$ units of gap and the window family for $116$.
Both figures are ceilings: a linear relaxation lands below them. The arc ceiling is the
one that matters, because it says the families now in the solvers have almost no
headroom left however they are tuned.}
\label{tab:ceilings}
\end{table}

Two readings follow. The arc-supported family is exhausted: on three quarters of the
loose instances the best bound it could ever certify is the coverage bound itself. And
the window family has real headroom, roughly two fifths of the remaining gap at the
median. Measured on the same instances, $84\%$ of the windows that carry the bound have
length at most four and $71\%$ have length two or three, so a static pool of short
windows captures most of the value and no separation routine is required.

\paragraph{Where the family belongs.}
In the position-indexed master, not the Benders master. Expressing \eqref{eq:window} in arc space needs a family of job
sets to be simultaneously contiguous, which is a product of indicators and has no useful
linear form --- the same obstruction that makes the triplet rows worthless. In the
position-indexed master the window is a range of position indices and the inequality is
immediate. That is an argument for the branch-and-price side of this report over the
Benders side, and it is the reverse of the ranking in Table~\ref{tab:solves}: the method
that solves fewer instances today is the one with somewhere left to go.

\subsection{Where the theory and the measurements meet}\label{ssec:disc-agree}

Three agreements between Section~\ref{sec:gap} and Section~\ref{sec:comp} are worth
drawing out, because each was arrived at independently.

\paragraph{The bound governs both sides.}
Section~\ref{sec:gap} shows that $q$ equals the optimum on a large part of the
instance space, and that where it does the two-phase heuristic tends to be exact.
Section~\ref{sec:comp} shows that where it does, the exact methods close at the root,
and where it does not they branch without closing. The same quantity that makes the
heuristic exact makes the exact methods fast. That is not obvious in advance: it says
that the instances on which one would most want an exact method are exactly the
instances on which the available exact methods are weakest.

\paragraph{The gap opens where the magazine is large.}
Observation~\ref{obs:census} finds that every instance with a positive heuristic gap
has $b\ge4$, and that positive gaps are rare in a random sample. Independently,
\citet{Burger2015} report from an industrial study that the decomposition does well
when the largest job nearly fills the magazine, and that its solution quality
decreases as the gap between the magazine size and the largest job size grows. One of
these is an exhaustive combinatorial search on small instances and the other is a case
study on a printing press. They describe the same regime.

\paragraph{The walk model overstates the weakness of the standard method.}
This is the least expected of the three. The literature's mental model of the
two-phase heuristic is a walk through one magazine per group, and under that model the
canonical example family has a positive gap that grows without bound. Under the method
as it is actually specified, with its final loading step, that same family has no gap
at all (Observation~\ref{obs:rings}). The standard method is better than the standard
picture of it. Problem~\ref{prob:ktns-closes} asks how much better, and that question
did not exist before the two readings were separated.

\subsection{Limitations}\label{ssec:limits}

The limitations are ordered by how much they threaten the report's conclusions, most
first, and each states the claim it bears on before the limit itself.

\paragraph{The census of heuristic gaps is on small instances.}
\emph{Bears on:} Observation~\ref{obs:census}, that positive heuristic gaps are rare,
equal to one, and require $b\ge4$.
\emph{The limit:} the census enumerates instances with three to five jobs, because it
computes the optimum by exhaustion. Nothing in it establishes that the same holds at
industrial sizes, and the agreement with \citet{Burger2015} is corroboration from a
different method, not an extension of the range.
\emph{Weight:} this is the least secure of the report's empirical claims, and the one a
reader should discount first. The structural results of Sections~\ref{ssec:exact} and
\ref{ssec:ratio} do not depend on it; they are proved for every instance and every
capacity.

\paragraph{The window family has been measured at its ceiling, not inside a solver.}
\emph{Bears on:} Table~\ref{tab:ceilings} and the claim that the family has roughly four
times the headroom of the arc-supported one.
\emph{The limit:} the figure is the strongest bound the family could certify, computed
by enumerating every sequence of $62$ instances of one family at eight jobs. A linear
relaxation will realise less than that, possibly much less, and no measurement inside a
branch-and-price run exists yet.
\emph{Weight:} the comparison between the two families is on equal terms, both being
ceilings, so the \emph{ratio} is meaningful even though neither absolute figure is
achievable. The claim to take from it is that the arc family is exhausted, which is a
ceiling result and therefore safe. The claim that the window family will deliver is not
yet established.

\paragraph{One time limit, one machine, one collection.}
\emph{Bears on:} every solve count and time in Section~\ref{sec:comp}.
\emph{The limit:} one hour, one homogeneous node type, and the standard collections. A
longer limit would move the counts; the cumulative curves of Table~\ref{tab:cumulative}
show by how much, and show that it would not move them equally for all methods.
\emph{Weight:} the conclusions of Section~\ref{ssec:disc-bound} rest on the split by
bound tightness and on the ablation, both of which are internal comparisons under
identical conditions and so are insensitive to where the limit is set.

\paragraph{The compact baseline is F4, and F4 is not the strongest of its family.}
\emph{Bears on:} the ranking in Table~\ref{tab:solves}.
\emph{The limit:} \citet{Catanzaro2015} recommend F5, which was not implemented.
\emph{Weight:} the omission runs against this report's own conclusion rather than for
it. A stronger compact baseline would widen a margin that already favours the compact
models over the methods built here. It has no bearing on the bound results, which
concern relaxation values rather than solver speed.

\paragraph{Four hundred and twenty-seven runs returned no result.}
\emph{Bears on:} the denominators in Table~\ref{tab:solves}.
\emph{The limit:} $352$ of the failures belong to the LSS baseline and $239$ of those to
one family, so the LSS row is computed on a smaller and harder-than-average subset and
its rate is not comparable with the others.
\emph{Weight:} the remaining $75$ failures are spread across nine configurations, at most
$26$ in any one, which changes no comparison drawn between them. The LSS row should be
read as a lower bound and nothing more.

\paragraph{The branch-and-price prototypes run in Python; the baselines run in C.}
\emph{Bears on:} the solve counts in Table~\ref{tab:bnp}.
\emph{The limit:} a node of a Python pricer costs orders of magnitude more than a node
inside CPLEX, and the run counts are additionally outcome-dependent through the
early-stop rule. Nothing above nine jobs was closed, and that fact carries no
information about the formulations.
\emph{Weight:} none of the conclusions drawn from these runs uses a solve count. They
rest on the root relaxation value, which is recorded before any search and is therefore
immune to both problems.

\paragraph{The learned cut-selection prototype was not run on the SSP.}
\emph{Bears on:} Section~\ref{ssec:rl} and nothing else.
\emph{The limit:} it was trained and measured on knapsack cover cuts, with one seed per
size and no repetition. It shows that the agent learns and says nothing about this
problem.
\emph{Weight:} Section~\ref{ssec:disc-locality} explains why running it on the SSP as
built would not have helped: cut selection chooses among cuts a method can already
generate, and the measurement there is that those cuts have almost no headroom left.

\paragraph{The setup-cost parameter has no measured range.}
\emph{Bears on:} the practical reading of Theorem~\ref{thm:collapse}, which locates the
transition between the SSP and the Job Grouping Problem at a threshold in the parameter
$\rho$, the ratio of the fixed cost of a machine stop to the cost of exchanging one
tool.
\emph{The limit:} no published source giving comparable measurements of the two
quantities was found, so where real installations sit on that spectrum is not
established here.
\emph{Weight:} the theorem is a statement about the model and is proved. What is missing
is an empirical value for one of its inputs, which affects how the result should be
applied and not whether it holds.

\paragraph{PTF exceeds the coverage bound, but only just.}
\emph{Bears on:} Proposition~\ref{prop:ptf} and the claim that it is the one construction
here to improve on the counting bound.
\emph{The limit:} the excess is one unit, on three runs of $233$, against gaps of six to
eight. The proposition says the relaxation \emph{can} exceed $q$ and is established by
witness; whether the excess can be made to grow is Problem~\ref{prob:ptf-gap} and is
open.
\emph{Weight:} the qualitative claim is secure and the quantitative one is not made.
